% Unit 10 teaching — Inverse Functions & King's/Queen's Rule
\input{technique-note.tex}
\title{Integration Bee --- Unit 10:\\ Inverse Functions \& King's/Queen's Rule}
\author{UK Uni Integration Bee prep\\ Syllabus reflection, Techniques Guide \S\S1.2 and 1.5}
\date{}
\begin{document}
\maketitle
\section*{Status}
Concept teaching. Tags: \texttt{inverse-functions}, \texttt{kings-rule}, \texttt{yang:1.17}, \texttt{yang:1.18}.
The Yang chapter is one source. The identity below is the syllabus reflection property and Guide \S1.2. The inverse-area step is Guide \S1.5. Nested and exponential copies are unit 25.
\section{Official anchors}
\begin{itemize}
  \item Syllabus \S1 --- reflection property $\int_a^b f(x)\,dx=\int_a^b f(a+b-x)\,dx$.
  \item Guide \S1.2 Reflection; \S1.5 Integral of an inverse.
  \item Yang \S\S1.17 Inverse Functions, 1.18 King's/Queen's Rule, as a later writeup of the same two moves.
\end{itemize}
\section{Learning goals}
\begin{enumerate}
  \item Integrate an inverse by parts, or by the area identity.
  \item Apply the one reflection $u=a+b-x$, including the sine-cosine case on $[0,\pi/2]$.
  \item Recognise the even, odd, and $x\mapsto 1/x$ cases as the same idea.
\end{enumerate}
\section{Pedagogy}
Worked $\to$ faded $\to$ interleaved. See \texttt{pedagogy/README.md}.

\section{Worked inverse --- \texttt{U10-D01} (\texttt{yang:1.17})}
\[
\int_0^1\arcsin x\,dx.
\]
Parts with $u=\arcsin x$, $dv=dx$ gives
\[
\bigl[x\arcsin x\bigr]_0^1 - \int_0^1\frac{x}{\sqrt{1-x^2}}\,dx.
\]
The boundary term is $\pi/2$. The remaining integral is $-\sqrt{1-x^2}$, so
\[
\frac\pi2 - \bigl[-\sqrt{1-x^2}\bigr]_0^1 = \frac\pi2 - (0-1) = \frac\pi2-1.
\]
The same value falls out of areas. On $[0,\pi/2]$, $\sin$ and $\arcsin$ are inverses, and
\[
\int_0^1\arcsin y\,dy = 1\cdot\frac\pi2 - \int_0^{\pi/2}\sin x\,dx = \frac\pi2-1.
\]
That area identity is the general one: if $f$ is increasing and continuous,
\[
\int_a^b f^{-1}(y)\,dy = b\,f^{-1}(b)-a\,f^{-1}(a)-\int_{f^{-1}(a)}^{f^{-1}(b)} f(x)\,dx,
\]
whenever the integrals exist.

\section{Faded inverse --- \texttt{U10-D08}}
\[
\int_0^{1/2}\arcsin x\,dx.
\]
The boundary term is $(1/2)\cdot(\pi/6)=\pi/12$. The same substitution $w=1-x^2$ produces $\sqrt3/2-1$ from the remaining integral. The stored value is $\pi/12+\sqrt3/2-1$.

\texttt{U10-D02} is $\arctan$ on $[0,1]$, parts again, stored value $\pi/4-(\ln 2)/2$. \texttt{U10-D03} is $\arccos$ on $[0,1]$, stored value $1$.

\section{Worked reflection}
The identity keeps the same interval. For a continuous $f$,
\[
\int_a^b f(x)\,dx=\int_a^b f(a+b-x)\,dx.
\]
The substitution $u=a+b-x$ swaps the limits and then the sign of $du$ swaps them back. Geometrically it reflects the graph in the line $x=(a+b)/2$, so the endpoints stay put. Contest writeups call this King's rule. They call the special case on $[0,a]$, namely $u=a-x$, Queen's rule. There is one identity.

If $f(x)+f(a+b-x)$ equals a constant $c$, the integral equals $c(b-a)/2$. Write both integrands before you add. If they do not add to a constant, keep both copies. The sum may still be shorter than either one.

Add the original integral to the reflected copy when the two integrands sum to something you can integrate. That is the whole toolkit.

On $[-a,a]$ the endpoints sum to zero, so $a+b-x$ is $-x$.
An even integrand is unchanged by that swap.
An odd integrand changes sign, so the integral is 0 when it exists.

On $[0,\pi/2]$ the useful reflection is $x\mapsto\pi/2-x$.
It swaps sine and cosine. For \texttt{U10-D07},
\[
I=\int_0^{\pi/2}\sin^2 x\,dx
=\int_0^{\pi/2}\cos^2 x\,dx,
\]
and $\sin^2 x+\cos^2 x=1$, so $2I=\pi/2$ and $I=\pi/4$.

\section{Worked King --- \texttt{U10-D04} (\texttt{yang:1.18})}
\[
I=\int_0^{\pi/2}\frac{\sin x}{\sin x+\cos x}\,dx.
\]
The swap $x\mapsto\pi/2-x$ exchanges the two summands, so the reflected integrand is $\cos x/(\sin x+\cos x)$. Adding the copies gives $2I=\int_0^{\pi/2} 1\,dx=\pi/2$, hence $I=\pi/4$.
\texttt{U10-D05} is the cosine-over-sum partner and has the same value. The King example \texttt{U10-D04} is the sine and cosine case, not the whole toolkit.

\section{Faded reflection --- \texttt{U10-D05}}
Cover the worked king. The integrand of \texttt{U10-D05} is $\cos x/(\sin x+\cos x)$ on $[0,\pi/2]$. The swap $x\mapsto\pi/2-x$ replaces it by a fraction. Write that fraction, then add. The two fractions add to 1, and the interval has length $\pi/2$. The stored value matches \texttt{U10-D04}.

The same complement, seen from a tangent, is still \texttt{U10-D07}. On $[0,\pi/2]$,
\[
\frac{1}{1+\tan^2 x}=\cos^2 x,
\]
and $\tan(\pi/2-x)=\cot x$, so the reflected fraction is $\tan^2 x/(1+\tan^2 x)$. The two add to 1, and half of the length $\pi/2$ is $\pi/4$. That is the stored value of \texttt{U10-D07}. You do not need a new drill for the tangent form.

\section{Log sine, not a drill}
The Techniques Guide evaluates $\int_0^{\pi/2}\ln(\sin x)\,dx$ by this reflection. The value is not a stored drill. Do not mark it as one.
Let $I$ be that integral. The swap $x\mapsto\pi/2-x$ gives $I=\int_0^{\pi/2}\ln(\cos x)\,dx$, so
\[
2I=\int_0^{\pi/2}\ln(\sin x\cos x)\,dx
=\int_0^{\pi/2}\ln(\sin 2x)\,dx-\frac\pi2\ln 2.
\]
The substitution $u=2x$ turns $\int_0^{\pi/2}\ln(\sin 2x)\,dx$ into $\frac12\int_0^\pi\ln(\sin u)\,du$. On $[\pi/2,\pi]$ the map $u\mapsto\pi-u$ copies $\ln(\sin u)$ onto $[0,\pi/2]$, so the integral from $0$ to $\pi$ is $2I$. Therefore
\[
2I=\frac12\cdot 2I-\frac\pi2\ln 2=I-\frac\pi2\ln 2,
\]
and
\[
I=-\frac\pi2\ln 2.
\]
A high-precision quadrature agrees with that value. The derivation above is the reason to quote it. SymPy leaves the symbolic integral unevaluated, so this is not a machine drill.

\section{Algebra on a reflected interval}
\texttt{U10-D06} is $\int_0^1 x/(1+x)\,dx$. Reflection sends it to $\int_0^1 (1-x)/(2-x)\,dx$, which is a different integral. The short path is the constant you add and subtract:
\[
\frac{x}{1+x}=1-\frac{1}{1+x},\qquad \int_0^1=1-\ln 2.
\]
That algebra is Guide \S1.9 and \texttt{docs/techniques/introduce-constant.md}. Use it when reflection produces a mess and the numerator is almost the denominator.

\section{The half-line}
On $(0,\infty)$ the substitution $x=1/u$ is the same idea at a different point.
It sends $\ln x/(1+x^2)$ to its negative, and $dx/(1+x^2)$ is unchanged, so
\[
\int_0^\infty\frac{\ln x}{1+x^2}\,dx=0
\]
when the integral exists. The stored drill is \texttt{U33-D07}. Unit 33 writes the general inversion; the sign flip is the reflection.

\section{Where the toolkit goes next}
Unit 25 uses the same $a+b-x$ on an exponential weight and on $x/(1+\sin x)$. If the reflected integrand is not the original and not a complement, you still add the two copies. The half-angle substitution is a different move. It lives in unit 09 and, as the named note, in unit 26.

\section{Practice}
\texttt{U10-D01}\ldots\texttt{U10-D08} in \texttt{drills.tex}. Machine check: \texttt{verify/run.sh}.

\section*{On a miss}
Teach this file. Rebuild the reflected copy and add it to the original. Do not invent an answer that is not in \texttt{drills.md}.
\end{document}
